% GENERATED FILE. Edit canonical lessons, then run npm run generate:books.
% canonical-source-hash: fnv1a64:696c7239873c9ddd
% canonical-lessons: MW-C219-padharno, MW-C219-bhet, MW-C219-haphto, MW-C219-tij, MW-C219-utavlo

\chapter{Guests Are Coming}
\label{ch:mw-guests-are-coming}
\hlchaptermodality{219}

\begin{chapteropening}
\textbf{By the end of this chapter:} \emph{I can ask a guest about coming and say what I will bring.}

Complete the last lesson of chapter 219: I can ask a guest about coming and say what I will bring.
\end{chapteropening}

\section[\texorpdfstring{\mw{पधारणो} (padhārṇo)}{पधारणो (padhārṇo)}]{\mw{पधारणो} (padhārṇo) — to come (said with respect)}

\label{lesson:MW-C219-padharno}



\begin{quote}
Before the new one: say the Marwadi for to depart, then the Marwadi for to take to a place.
\end{quote}

\subsection*{You'll want to know: \mw{पधारणो}}
\textbf{\mw{पधारणो}} — \emph{padhārṇo} — \textquotedblleft{}to come (said with respect)\textquotedblright{}.

\begin{quote}
\mw{आप} \mw{काल} \mw{पधारो} \mw{हो}?
\end{quote}

\begin{quote}
\emph{āp kāl padhāro ho?}
\end{quote}

\begin{quote}
Will you be coming tomorrow?
\end{quote}

Past (he …): \mw{वो} \mw{घरै} \mw{पधार्यो}, \emph{vo gharai padhāryo}.

\begin{cousinweb}
\textbf{\mw{पधारणो}} (\emph{padhārṇo}) is inherited from a reconstructed \emph{*paddhārayati}, \textquotedblleft{}steps\textquotedblright{}; CDIAL lists the Old Marwari verb \emph{padhāraï}, \textquotedblleft{}comes, goes\textquotedblright{} (CDIAL 7768). It is the polite word for a guest's coming.
\end{cousinweb}

\begin{grammarlens}[title={What you've built}]
One more word for saying what comes next.
\end{grammarlens}

\subsection*{Your turn}
\begin{itemize}
\raggedright
  \item \emph{Say it:} \emph{padhārṇo}
  \item \emph{Say it:} \emph{padhārṇo}, once more
  \item \emph{Say it:} \emph{āp kāl padhāro ho?} — Will you be coming tomorrow?
  \item \emph{Recall:} say the Marwadi for to depart, then the Marwadi for to take to a place, then say \emph{padhārṇo} again
\end{itemize}

\begin{tcolorbox}[breakable,colback=teal!4,colframe=teal!35!black,title={Before you move on}]
What does \mw{पधारणो} mean? (\textquotedblleft{}To come (said with respect)\textquotedblright{}.) Say it once more.
\end{tcolorbox}
\section[\texorpdfstring{\mw{भेट} (bheṭ)}{भेट (bheṭ)}]{\mw{भेट} (bheṭ) — a gift (brought on a visit); a meeting}

\label{lesson:MW-C219-bhet}



\begin{quote}
Before the new one: say the Marwadi for to take to a place, then the Marwadi for to come (said with respect).
\end{quote}

\subsection*{You'll want to know: \mw{भेट}}
\textbf{\mw{भेट}} — \emph{bheṭ} — \textquotedblleft{}a gift (brought on a visit); a meeting\textquotedblright{}.

\begin{quote}
\mw{हूं} \mw{काल} \mw{भेट} \mw{लासूं।}
\end{quote}

\begin{quote}
\emph{hū̃ kāl bheṭ lāsū̃.}
\end{quote}

\begin{quote}
Tomorrow I'll bring a gift.
\end{quote}

A feminine noun.

\begin{cousinweb}
\textbf{\mw{भेट}} (\emph{bheṭ}) is inherited from a reconstructed \emph{*bhiṭ}, \textquotedblleft{}meet\textquotedblright{}; CDIAL lists Old Marwari \emph{bheṭi}, \textquotedblleft{}a meeting\textquotedblright{}, and Hindi \emph{bheṭī} and Sindhi \emph{bheṭa}, \textquotedblleft{}a present\textquotedblright{} (CDIAL 9490): the gift you bring to a meeting.
\end{cousinweb}

\begin{grammarlens}[title={What you've built}]
One more word for saying what comes next.
\end{grammarlens}

\subsection*{Your turn}
\begin{itemize}
\raggedright
  \item \emph{Say it:} \emph{bheṭ}
  \item \emph{Say it:} \emph{bheṭ}, once more
  \item \emph{Say it:} \emph{hū̃ kāl bheṭ lāsū̃.} — Tomorrow I'll bring a gift.
  \item \emph{Recall:} say the Marwadi for to take to a place, then the Marwadi for to come (said with respect), then say \emph{bheṭ} again
\end{itemize}

\begin{tcolorbox}[breakable,colback=teal!4,colframe=teal!35!black,title={Before you move on}]
What does \mw{भेट} mean? (\textquotedblleft{}A gift (brought on a visit); a meeting\textquotedblright{}.) Say it once more.
\end{tcolorbox}
\section[\texorpdfstring{\mw{हफ्तो} (haphto)}{हफ्तो (haphto)}]{\mw{हफ्तो} (haphto) — a week}

\label{lesson:MW-C219-haphto}



\begin{quote}
Before the new one: say the Marwadi for to come (said with respect), then the Marwadi for a gift.
\end{quote}

\subsection*{You'll want to know: \mw{हफ्तो}}
\textbf{\mw{हफ्तो}} — \emph{haphto} — \textquotedblleft{}a week\textquotedblright{}.

\begin{quote}
\mw{हूं} \mw{एक} \mw{हफ्तो} \mw{गांव} \mw{मांय} \mw{रैसूं।}
\end{quote}

\begin{quote}
\emph{hū̃ ek haphto gā̃v mā̃y raisū̃.}
\end{quote}

\begin{quote}
I'll stay a week in the village.
\end{quote}

A masculine noun.

\begin{cousinweb}
\textbf{\mw{हफ्तो}} (\emph{haphto}) is usually traced to Persian \emph{hafta}, \textquotedblleft{}a week\textquotedblright{}, from \emph{haft}, \textquotedblleft{}seven\textquotedblright{}; Marwadi ends it in \emph{-o} like its own nouns.
\end{cousinweb}

\begin{grammarlens}[title={What you've built}]
One more word for saying what comes next.
\end{grammarlens}

\subsection*{Your turn}
\begin{itemize}
\raggedright
  \item \emph{Say it:} \emph{haphto}
  \item \emph{Say it:} \emph{haphto}, once more
  \item \emph{Say it:} \emph{hū̃ ek haphto gā̃v mā̃y raisū̃.} — I'll stay a week in the village.
  \item \emph{Recall:} say the Marwadi for to come (said with respect), then the Marwadi for a gift, then say \emph{haphto} again
\end{itemize}

\begin{tcolorbox}[breakable,colback=teal!4,colframe=teal!35!black,title={Before you move on}]
What does \mw{हफ्तो} mean? (\textquotedblleft{}A week\textquotedblright{}.) Say it once more.
\end{tcolorbox}
\section[\texorpdfstring{\mw{तीज} (tīj)}{तीज (tīj)}]{\mw{तीज} (tīj) — the Teej festival; the third day of a lunar fortnight}

\label{lesson:MW-C219-tij}



\begin{quote}
Before the new one: say the Marwadi for a gift, then the Marwadi for a week.
\end{quote}

\subsection*{You'll want to know: \mw{तीज}}
\textbf{\mw{तीज}} — \emph{tīj} — \textquotedblleft{}the Teej festival; the third day of a lunar fortnight\textquotedblright{}.

\begin{quote}
\mw{परसूं} \mw{तीज} \mw{है।}
\end{quote}

\begin{quote}
\emph{parsū̃ tīj hai.}
\end{quote}

\begin{quote}
The day after tomorrow is Teej.
\end{quote}

A feminine noun.

\begin{cousinweb}
\textbf{\mw{तीज}} (\emph{tīj}) is inherited from Sanskrit \emph{tṛtīyā}, \textquotedblleft{}the third day of the lunar fortnight\textquotedblright{} (CDIAL 5920, which lists Marwari \emph{tīj}). The women's festival of Teej falls on such a third day, in the rains.
\end{cousinweb}

\begin{grammarlens}[title={What you've built}]
One more word for saying what comes next.
\end{grammarlens}

\subsection*{Your turn}
\begin{itemize}
\raggedright
  \item \emph{Say it:} \emph{tīj}
  \item \emph{Say it:} \emph{tīj}, once more
  \item \emph{Say it:} \emph{parsū̃ tīj hai.} — The day after tomorrow is Teej.
  \item \emph{Recall:} say the Marwadi for a gift, then the Marwadi for a week, then say \emph{tīj} again
\end{itemize}

\begin{tcolorbox}[breakable,colback=teal!4,colframe=teal!35!black,title={Before you move on}]
What does \mw{तीज} mean? (\textquotedblleft{}The Teej festival; the third day of a lunar fortnight\textquotedblright{}.) Say it once more.
\end{tcolorbox}
\section[\texorpdfstring{\mw{उतावळो} (utāvḷo)}{उतावळो (utāvḷo)}]{\mw{उतावळो} (utāvḷo) — in a hurry; hasty}

\label{lesson:MW-C219-utavlo}



\begin{quote}
Before the new one: say the Marwadi for a week, then the Marwadi for the Teej festival.
\end{quote}

\subsection*{You'll want to know: \mw{उतावळो}}
\textbf{\mw{उतावळो}} — \emph{utāvḷo} — \textquotedblleft{}in a hurry; hasty\textquotedblright{}.

\begin{quote}
\mw{हूं} \mw{उतावळो} \mw{हूं} \mw{क्यूंकै} \mw{बस} \mw{आथण} \mw{छूटै} \mw{है।}
\end{quote}

\begin{quote}
\emph{hū̃ utāvḷo hū̃ kyū̃kai bas āthaṇ chhūṭai hai.}
\end{quote}

\begin{quote}
I'm in a hurry because the bus leaves in the evening.
\end{quote}

\begin{cousinweb}
\textbf{\mw{उतावळो}} (\emph{utāvḷo}) goes back to a reconstructed \emph{*uttāpala}, \textquotedblleft{}impetuous\textquotedblright{}, built on Sanskrit \emph{uttāpa}, \textquotedblleft{}heat\textquotedblright{}; CDIAL lists a Gujarati word for impetuous under it. CDIAL gives no Marwari form in this entry (CDIAL 1788).
\end{cousinweb}

\begin{grammarlens}[title={What you've built}]
One more word for saying what comes next.
\end{grammarlens}

\subsection*{Your turn}
\begin{itemize}
\raggedright
  \item \emph{Say it:} \emph{utāvḷo}
  \item \emph{Say it:} \emph{utāvḷo}, once more
  \item \emph{Say it:} \emph{hū̃ utāvḷo hū̃ kyū̃kai bas āthaṇ chhūṭai hai.} — I'm in a hurry because the bus leaves in the evening.
  \item \emph{Recall:} say the Marwadi for a week, then the Marwadi for the Teej festival, then say \emph{utāvḷo} again
\end{itemize}

\begin{tcolorbox}[breakable,colback=teal!4,colframe=teal!35!black,title={Before you move on}]
What does \mw{उतावळो} mean? (\textquotedblleft{}In a hurry; hasty\textquotedblright{}.) Say it once more.
\end{tcolorbox}
